%======================================================================
\section{들어가며}
\label{sec:intro}

\begin{yoji}
블록체인 평판 점수는 보통 송금 그래프를 걸어요. 그런데 같은 주소들은 허락으로도 이어져 있어요. 이 원고는 두 관계를 한 걷기로 섞으면 계약의 순위를 더 잘 매기는지를, 미리 정해 커밋한 두 규칙으로 시험해요. 두 규칙은 모두 실패했어요.
\end{yoji}

\subsection{디파이는 주소로만 아는 상대와 거래해요}

디파이(DeFi, 탈중앙 금융)의 대출과 마진 거래 프로토콜은 주소로만 아는 상대와 연동을 다뤄요~\cite{schar2021defi}. 그래서 신용 기록 대신 과담보(빌리는 돈보다 담보를 더 맡기기)와 자동 청산에 기대요~\cite{gudgeon2020crisis,werner2022sok}. 한 갈래의 연구는 빠진 정보를 공개 장부에서 얼마쯤 되찾으려 해요. 그 대부분은 토큰 송금 그래프 위에서 주소의 순위를 매겨요.
\begin{itemize}
\item Do와 Do는 이더리움 거래 위의 페이지랭크와 호지랭크(HodgeRank)를 사회적 신용의 척도로 읽어요~\cite{do2023hodgerank}.
\item Do, Do, Nguyen의 AWP(적응형 가중 페이지랭크)는 송금마다 금액과 나이로 무게를 달고, 최근에 활발히 송금한 주소에서 걷기를 다시 출발시켜요~\cite{do2023awp}.
\item RiskProp은 거래 네트워크를 따라 계정의 위험을 퍼뜨려요~\cite{lin2025riskprop}.
\end{itemize}
블록체인 평판 시스템을 살핀 조사도 같은 의존, 곧 송금량이나 명시적 평가에 기대는 모습을 찾아요~\cite{bellini2020trust}. 이런 점수는 누구나 공개 사건으로 다시 계산할 수 있어요. 그래서 탈중앙 신용 평가가 속을 알 수 없는 블랙박스가 된다는 걱정에 얼마쯤 답해요~\cite{packin2024credit}.

\subsection{장부는 두 번째 관계도 기록해요}

송금은 지급을 기록해요. 그런데 장부는 같은 주소들 사이의 두 번째 관계도 기록해요. ERC-20(Ethereum Request for Comments 20) 토큰 규칙에서 주인은 \texttt{approve}를 불러, 스펜더가 자기 토큰을 정한 양까지 꺼내 쓰게 해요. EIP 2612(Ethereum Improvement Proposal)의 \texttt{permit} 함수는 블록체인 밖의 서명으로 같은 허락을 정해요~\cite{vogelsteller2015erc20,lundfall2020permit}.

허락은 돈을 옮기지 않아요. 스펜더가 나중에 쓸 수 있는 능력을 줄 뿐이고, 주인은 그 허락을 올리거나 내리거나 거둘 수 있어요. 그리고 많은 허락이 무제한이에요~\cite{wang2022unlimited}. 허락은 주로 사용자 계정에서 그 계정이 쓰는 라우터, 금고(vault), 거래 계약 쪽을 가리켜요. 그래서 허락은 스마트 계약의 순위를 매기는 데 자연스러운 신호예요. 이렇게 지급(송금)과 허락은 한 주소 집합 위의 두 층을 이뤄요. 다층 순위 매기기가 다루는 바로 그 상황이에요~\cite{halu2013multiplex,dedomenico2015ranking}.

\subsection{이 원고가 묻는 것}

이 원고는 이렇게 물어요. 두 층을 함께 걷는 무작위 걷기 하나가, 한 층만 걷는 것보다 계약의 순위를 더 잘 매길까요? 특히 송금 걷기에 허락 화살표를 더하면 나아질까요? 이를 위해 두 가지를 정의해요.
\begin{itemize}
\item \textbf{C-PR(결합 페이지랭크):} 점마다 두 층의 이동 행을 섞어요. 허락 층에는 고정된 무게 $\lambda$를 줘요.
\item \textbf{S-PR(보증 씨앗 페이지랭크):} 빼 보기 실험이에요. 송금 층을 걷되, 다시 출발할 곳을 허락 걷기에서 뽑아요. 트러스트랭크(TrustRank)가 미리 검증한 씨앗 집합에서 다시 출발하는 것과 같은 생각이에요~\cite{gyongyi2004trustrank}.
\end{itemize}
C-PR 가족의 끝점인 $\lambda=1$과 $\lambda=0$은 한 층만 걸어요. 이 둘이 두 판정 규칙의 비교 대상이에요. 점수를 받는 것은 처음부터 끝까지 스마트 계약이에요. 외부 소유 계정은 주인, 보낸 쪽, 받는 쪽으로 그래프에 들어가지만, 점수를 받지 않아요.

\subsection{시간을 앞으로 보는 평가}

평가는 시간을 앞으로 봐요. 어느 날짜에 점수를 동결하고, 그 뒤에 일어나는 결과와의 켄달 순위 상관으로 판단해요~\cite{kendall1938rank}. 결과는 세 가지예요.
\begin{itemize}
\item 그 계약에 새로 허락하는 주인.
\item 그 계약에 새로 송금하는 주소.
\item GMX V2 무기한 선물 거래소~\cite{gmxnddocs}에서 거래하는 계약이라면, 포지션 닫기 가운데 강제 청산을 피한 몫.
\end{itemize}
앞의 둘은 층을 이루는 것과 같은 종류의 화살표가 새로 생기는 것이에요. 셋째는 어느 층으로도 만들지 않았고, 이 자료에서 상대방 위험에 가장 가까운 결과예요~\cite{perez2021liquidations}.

\subsection{미리 정한 두 규칙과 그 결과}

판정 규칙 두 개를, 이 연구가 그 규칙으로 판정할 자료를 추출하기 전에 정했어요.
\begin{itemize}
\item \textbf{분석 계획:} 규칙 A와 모든 점수를 정했어요. 관측 기간을 복사한 BigQuery 작업이 만들어지기 10초 전에 저자들의 공개 저장소에 커밋했어요.
\item \textbf{두 번째 창의 등록:} 비열등성 허용 폭이 있는 규칙 B를 정했어요. 2026년 7–9월 로그를 복사한 작업보다 8초 먼저 커밋했어요. 그 작업을 돌리는 스크립트는 등록 파일이 로컬 수정 없이 커밋되어 있지 않으면 시작하지 않아요. 시작하면 그 파일들의 SHA-256 해시를 매니페스트에 적어요.
\end{itemize}
이 원고는 이런 방식을 “커밋으로 잠근(commit-gated)” 방식이라고 불러요. 두 커밋이 추출보다 먼저였다는 것은 커밋 시각과 작업 시각에 기대요. 제3자의 시간 도장에 기대는 것은 아니에요(\ref{sec:gate}절).

두 규칙은 모두 실패했어요. 이 결과를 안 뒤, 저자들은 계획이 금액에 무게를 다는 방식에서 결함 세 가지를 찾았어요. 그리고 이 원고가 먼저 보고하는 분석을 다음 세 가지로 정의했어요.
\begin{itemize}
\item 가격이 믿을 만한 다섯 토큰의 달러 금액.
\item 무제한 허락에 두는 상한.
\item 동결 직전 12개월의 사건으로 만든 점수.
\end{itemize}
계획대로라면 금액이 아무 효과가 없거나, 무제한 허락이 행을 다 차지했어요. 그래서 이 원고는 이 수정 분석을 앞에 둬요. 하지만 그 정의(달러 금액, 다섯 토큰, 12개월 점수 기간)가 결과 뒤에 나왔으므로, 그 판정은 같은 규칙에 따른 재분석이에요. 미리 정한 분석은 같은 판정에 이르렀고, 결과의 모든 표와 그림에서 수정 분석 옆에 함께 보고해요.

\subsection{기여}

\begin{enumerate}
\item \textbf{계약 평판의 두 층 모형.} 결합 걷기, 씨앗을 쓰는 빼 보기 실험, 분석적 관찰 두 가지를 내놓아요.
  \begin{itemize}
  \item $\lambda$가 0과 1 사이에 있을 때, 무게는 두 층 모두에 나가는 화살표가 있는 점을 통해서만 작용해요.
  \item 스펜더를 허락한 주인들이 스스로는 허락을 받지 않는 그래프에서, 허락 걷기는 주인 쪽에서 정규화한 받은 무게로 줄어들어요.
  \end{itemize}
\item \textbf{장부 자료 위의 순위 실험을 위한, 미리 정하고 커밋으로 잠근 절차.} 약속마다 날짜를 매기는 기록(\ref{sec:protocol}절)과 그 기록의 한계를 함께 내놓아요. 커밋 시각은 커밋한 사람의 시계가 정하고, 제3자의 시간 도장이 그 시각을 붙잡아 두지 않아요.
\item \textbf{아비트럼 원에서의 부정적 결과.} C-PR은 수정 분석에서도, 미리 정한 분석에서도 어느 규칙도 만족하지 못했어요.
  \begin{itemize}
  \item 나중의 허락에서 송금 걷기보다 나은 점은, 규칙을 정한 뒤 이 연구가 추출한 창에서도 유지됐어요.
  \item 하지만 나중의 송금자에서는 등록한 허용 폭보다 더 크게 뒤졌고, AWP에는 그보다 더 크게 뒤졌어요.
  \item 트레이더 결과에서는 허용 폭 안에 머물지 못했어요.
  \item 라벨을 안 뒤 계산한 짝지은 대비에서, C-PR은 수정 분석의 두 창 모두에서 나중의 허락은 허락한 주인 수만큼, 나중의 송금자는 서로 다른 송금자 수만큼 잘 줄 세우지 못했어요.
  \end{itemize}
\item \textbf{알려진 링크 농장 결과를 균등 재시작 아래의 허락 층에 맞춰 잰 사후 보정.} 수정 분석의 엔도스랭크 그래프에서, 닫힌 농장은 감쇠 계수 0.85에서 자기 인구 몫의 3.23배를 모아요. 농장 주소 26개, 약 2,292개, 약 47,073개짜리 별 모양은 목표를 코호트 엔도스랭크 점수의 90번째 백분위, 99번째 백분위, 꼭대기 위로 올려요. 등록 분석의 그래프에서는 이 수가 2.96, 13, 1,253, 약 1,080,271이에요.
\end{enumerate}

\subsection{짝 원고와 나눠 맡은 것}

짝 원고~\cite{kwon2026companion}는 두 화살표 종류를 따로따로의 점수로 견주고, 각 점수를 그 원시 차수와 견줘요. 그리고 어느 화살표로도 만들지 않은 트레이더 결과를 세 규칙(R1–R3)으로 다뤄요. 이 세 규칙은 규칙 B와 같은 커밋에 등록했어요.
\begin{itemize}
\item \textbf{두 원고가 함께 보고하는 것:} 같은 자료, 스펜더 홀드아웃에서 원시 차수·엔도스랭크·AWP의 계수와 차수 대비, 허락 걷기의 닫힌 식, 풀이 실행 시간.
\item \textbf{짝 원고에만 있는 것:} 트레이더 결과에서 허락 쪽과 송금 쪽을 견준 비교(그 연구의 규칙 R1–R3).
\item \textbf{이 원고가 같은 트레이더 코호트에서 보고하는 것:} 규칙 B와 그 등록 민감도 분석이 이름 붙인 C-PR의 대비, 그리고 라벨을 안 뒤에 계산한 C-PR과 받은 허락 수의 대비 하나.
\end{itemize}
두 원고 모두 수정 분석을 먼저 보고하고 등록 분석의 값을 옆에 둬요. 그리고 두 원고는 두 분석 모두에서 서로 같은 값을 보고해요. 이 원고가 다루는 것은 결합 연산자, 미리 정한 두 규칙과 그 기록, 그리고 허락 층을 얼마나 조작할 수 있는지예요.

글의 차례는 이래요. \ref{sec:related}절은 관련 연구, \ref{sec:model}절은 모형, \ref{sec:protocol}절은 규칙과 그 기록, \ref{sec:data}절은 자료와 측정, \ref{sec:results}절은 결과, \ref{sec:sybil}절은 농장 공격의 보정, \ref{sec:threats}절은 타당성을 위협하는 것, \ref{sec:implications}절은 대출·마진 프로토콜에 주는 시사점, \ref{sec:conclusion}절은 결론이에요. 부록~\ref{app:explore}에는 탐색적 점검을 모았어요.

%======================================================================
\section{관련 연구}
\label{sec:related}

\begin{yoji}
장부 그래프의 평판 연구는 거의 송금만 걸었어요. 다층 순위 매기기와 개인화 페이지랭크는 두 층을 섞는 방법을 줘요. 차수는 링크 예측의 표준 기준선이라, 걷기는 그 걷기가 다듬는 원시 차수와 견줘 시험해야 해요. 페이지랭크 같은 대칭 함수는 시빌 공격에 안전할 수 없어요. 미리 정한 분석은, 분석하는 사람이 아직 결과를 살피지 않았다면 이미 있는 자료에도 쓸 수 있어요.
\end{yoji}

\subsection{장부 그래프로 만든 평판}

페이지랭크는 방향이 있는 그래프의 점들을, 감쇠 계수를 쓰는 무작위 걷기의 정상 분포로 줄 세워요~\cite{page1999pagerank,langville2006pagerank}. HITS(Hyperlink-Induced Topic Search)는 허브와 권위를 나누고~\cite{kleinberg1999hits}, 아이겐트러스트(EigenTrust)는 지역 신뢰 값을 정규화해서 파일 공유 네트워크의 참여자가 자기 위치를 부풀리지 못하게 해요~\cite{kamvar2003eigentrust}.

공개 장부에서는 이런 생각을 거의 송금에만 적용했어요. 위에서 말한 방법들~\cite{do2023hodgerank,do2023awp} 말고도 이런 연구가 있어요. 거래 그래프의 네트워크 임베딩은 계정마다의 특징보다 피싱 계정을 더 잘 찾아요~\cite{wu2022phishers}. 암호화폐 거래의 지식 발견을 살핀 조사는 차수와 금액에 바탕을 둔 중심성을 흔한 주소 서술자로 꼽아요~\cite{liu2021knowledge}.

허락은 문헌에 두 가지 역할로 나와요. 한 주체의 주소들을 묶는 휴리스틱으로~\cite{victor2020clustering}, 그리고 무제한일 때의 보안 위험으로요~\cite{wang2022unlimited}. 휴먼 패스포트(Human Passport) 같은 도구는 무게를 단 자격 증명으로 주소의 주인을 점수 매길 뿐, 주소 사이 관계의 순위를 매기지는 않아요~\cite{humanpassportnddocs}. 짝 연구~\cite{kwon2026companion}는 마지막 허락을 순위 매기기의 화살표로 써요. 이 원고는 그 화살표를 송금 화살표와 한 걷기 안에서 묶어요.

\subsection{다층 순위와 개인화 순위}

\begin{itemize}
\item 멀티플렉스 페이지랭크는 한 층에서 얻은 중심성이 다른 층의 걷기를 기울이게 해요~\cite{halu2013multiplex}.
\item De Domenico 외의 서로 이어진 다층 틀에서는, 걷는 사람이 층 안에서 움직이다가 층 사이를 갈아타요. 그리고 점들을 여러 층에 걸친 쓰임새(versatility)로 줄 세워요~\cite{dedomenico2015ranking}.
\item 주제별 페이지랭크와 개인화 페이지랭크는 이동 행렬 대신 순간 이동 벡터를 바꿔요~\cite{haveliwala2003topic,jeh2003personalized}.
\item 트러스트랭크는 손으로 확인한 씨앗 집합으로만 순간 이동하게 해서 웹 스팸을 아래로 내려요~\cite{gyongyi2004trustrank}.
\end{itemize}
C-PR은 앞의 다층 걷기들(멀티플렉스 페이지랭크, 서로 이어진 다층 틀)의 특수한 경우예요. 행의 합이 1인 행렬 두 개를 점마다 볼록 결합(가중 평균)하되, 무게는 전체에 하나뿐이고 층 사이에 점의 복사본을 두지 않아요. S-PR은 순간 이동 벡터를 바꾸는 방법들(주제별·개인화 페이지랭크, 트러스트랭크)에 속해요. 아래의 명제~\ref{prop:localize}은 페이지랭크에서 흔히 쓰는 1차 섭동 한계~\cite{langville2006pagerank}를 이 볼록 결합에 맞춘 것이에요. 같은 주소들 사이의 두 온체인 관계에서, 연구가 평가 자료를 추출하기 전에 정한 규칙으로 이런 구성을 시험한 연구를 저자들은 문헌에서 찾지 못했어요.

\subsection{차수 기준선}

선호적 연결에서는 새 링크가 이미 링크가 많은 점에 몰려요~\cite{barabasi1999scaling}. 자라는 네트워크를 잰 연구도 같은 경향을 보여요~\cite{newman2001clustering}. 그래서 차수에 바탕을 둔 점수는 링크 예측의 표준 기준선이고~\cite{libennowell2007link}, 이더리움 거래 네트워크에서도 그래요~\cite{lin2022tracking}. 그리고 차수끼리의 상관이 약한 그래프에서는 페이지랭크가 평균적으로 들어오는 화살표 수에 비례해요~\cite{fortunato2008indegree}. 따라서 걷기가 정보를 더한다는 주장은, 그 걷기가 다듬는 원시 차수와 견줘 시험해야 해요. 이 원고는 모든 걷기를 그것과 견줘요.

\subsection{조작과 시빌 저항}

신원이 공짜인 곳에서는 한 행위자가 여럿으로 보일 수 있어요~\cite{douceur2002sybil}. Gy\"{o}ngyi와 Garcia-Molina는 목표의 페이지랭크를 가장 많이 올리는 링크 농장과 동맹을 끌어냈어요~\cite{gyongyi2005alliances}. Cheng과 Friedman은 페이지랭크를 포함해 어떤 대칭 평판 함수도 시빌 공격에 안전할 수 없다는 것을 보였어요. 반면 믿을 만한 출처에서 나가는 흐름이나 경로로 계산하는 비대칭 함수 가운데는 안전한 것도 있어요~\cite{cheng2005sybilproof}. \ref{sec:sybil}절은 이 결과를 허락 층에 적용해요. 새 결과를 끌어내지 않고, 수치로 보정할 뿐이에요.

장부에서는 자전 거래가 송금 활동을 부풀려요~\cite{victor2021washtrading,cong2023washtrading}. 사람 증명 방식은 시빌 저항과 개인정보 보호를 맞바꿔요~\cite{siddarth2020watchmen}. 그리고 탈중앙 시스템을 살핀 최근 검토는 자격 증명 파밍을 그런 관문의 남은 약점으로 꼽아요~\cite{bordeianu2026sybil}.

\subsection{미리 정한 분석}

사전 등록은 가설을 시험하는 분석과 가설을 만드는 분석을 나눠요. 분석하는 사람이 아직 결과를 살피지 않았다면, 이미 있는 자료에도 사전 등록을 쓸 수 있어요~\cite{nosek2018preregistration}. 관측된 장부 자료에서는 시험 창의 사건이 일어나는 순간부터 체인 위에 있어요. 그래서 연구가 줄 수 있는 가장 강한 현실적 보증은, 규칙을 정할 때 그 사건을 아직 추출하지 않았다는 것이에요. 구성 타당도 이론은 뒤에서 중요해지는 주의를 하나 더해요. 점수는 그것이 따라가도록 설계된 것의 점검과 맞아야 하므로, 그런 일치는 독립적인 증거가 아니에요~\cite{cronbach1955construct}.
